\section{Introduction}
\label{sec:intro}

Personalized neoantigen vaccines have moved from concept to the clinic, with
early-phase trials reporting durable, mutation-specific T-cell responses and
encouraging signals of clinical benefit across several tumour types
\cite{Sahin2023,NatRevImmunol2023}. Their central computational bottleneck is
selection: from the hundreds of somatic mutations in a tumour, only a handful of
mutated peptides will actually be recognized by a patient's T cells, and only a
fraction of those will mount a response strong enough to matter clinically.
Decades of work have made it clear that peptide--MHC binding, while necessary, is
far from sufficient for immunogenicity. The TESLA consortium benchmark made this
gap quantitative: of the strongly presented candidate peptides nominated across
participating teams, only on the order of six percent proved experimentally
immunogenic \cite{TESLA2020}. Prioritizing neoantigens for a vaccine therefore
demands predictors that reach beyond binding into the immunological consequences
of presentation.

The literature has approached this problem along what is best described as a
capability ladder. At the first level (L1), binding predictors such as
NetMHCpan and MHCflurry estimate peptide--MHC affinity or eluted-ligand
likelihood across thousands of alleles \cite{NetMHCpan2020,MHCflurry2020}. At
the second level (L2), presentation-aware models such as PRIME and BigMHC
incorporate mass-spectrometry eluted-ligand evidence to estimate the probability
that a peptide is genuinely presented \cite{PRIME2023,BigMHC2023}. At the third
level (L3), a large and architecturally diverse family of immunogenicity tools
asks the binary question of whether a presented peptide will be recognized by T
cells at all
\cite{DeepImmuno2021,deepHLApan2019,IMPROVE2024,PredIG2025,ICERFIRE2024,NeoaPred2024,ImmunoStruct2025,NetTepi2014,Seq2Neo2022,TSCAPE2025,CNNeoPP2026}.
A fourth concern, the \emph{magnitude} of the elicited response---how strong the
T-cell reaction is, rather than merely whether it occurs---remains essentially
unaddressed. We stress that this fourth concern is not a strictly higher rung of
the same ladder. Levels L1 through L3 sharpen a single binary judgement, whereas
magnitude is a \emph{regression axis orthogonal} to that judgement: one can ask
how strong a response will be at the presentation layer just as well as at the
recognition layer. We therefore treat magnitude (L4) as an orthogonal
quantitative dimension layered across the binary ladder, not as a successor to
it.

This orthogonal magnitude dimension is the gap that motivates our study.
Surveying more than a dozen immunogenicity methods published between 2024 and
2026, together with six recent reviews, we find that no tool is explicitly
designed to regress the continuous magnitude of the T-cell response and report
correlation or error against a continuous experimental readout; every method
ultimately emits a binary call or a binding-strength score
\cite{explorationpub2023,Zhao2025,BiomarkerRes2025}. The gap is striking
precisely because the signal is often available and then discarded: PredIG is
trained on data carrying graded positive-low, intermediate, and high labels that
are binarized before use \cite{PredIG2025}; ICERFIRE explicitly collapses graded
labels into a single class \cite{ICERFIRE2024}; and CNNeoPP's validation ELISpot
data distinguish weak from strong responders yet the model still outputs only a
binary decision \cite{CNNeoPP2026}. That a gap has gone unfilled, however, is not
by itself evidence that filling it is worthwhile. An absence in the literature
admits two readings: an open opportunity, or a problem that practitioners have
quietly judged intractable. Sequence-level magnitude prediction has a genuine
biological ceiling---the magnitude of a polyclonal T-cell response depends
heavily on donor-specific precursor frequencies that no peptide-and-allele input
can observe, plausibly capping the explainable variance at a Spearman
correlation of roughly $0.4$ to $0.6$. We take this ceiling seriously throughout,
and frame the value of continuous prediction not as high-precision regression but
as recovering ranking-relevant signal for clinical top-$K$ prioritization that
binary tools throw away.

Before a new magnitude predictor can be justified, however, a prior question must
be answered empirically: how well do today's tools actually track response
magnitude when their scores are repurposed as continuous predictors, and are the
apparent differences between them real? This question has not been asked under a
controlled, apples-to-apples protocol, in part because immunogenicity tools
differ in input formats, allele handling, score scales, and aggregation
conventions, so that head-to-head numbers are rarely comparable across papers. We
close this gap with a unified benchmark. Using a common ELISpot ground truth and
a seven-step harmonization protocol that standardizes inputs, alleles, and
scoring across tools, we evaluate eight immunogenicity tools and one presentation
proxy against the continuous ELISpot spot-forming-cell readout, and against the
binary positive/negative outcome, on a common patient cohort.

Our findings and contributions are threefold.
\textbf{First}, under this harmonized protocol the correlation between existing
tool scores and ELISpot magnitude is uniformly weak or non-significant: all tools
fall below an absolute Spearman correlation of $0.33$, with only IMPROVE and
PredIG reaching significance under particular aggregation choices, and the
binary-discrimination performance, while higher in point estimate, has bootstrap
confidence intervals that overlap so broadly---and shift so much with aggregation
and threshold choices---that no single tool can be declared best. A presentation
proxy performs near chance on magnitude, reinforcing that presentation,
immunogenicity, and response magnitude are distinct quantities. \textbf{Second},
we show that the global correlation metric used by default in the field masks
substantial individual variation, and we introduce a per-patient evaluation
protocol---computing a correlation within each patient and then aggregating---that
exposes it: tools whose pooled global correlation is indistinguishable from noise
recover meaningful within-patient ranking ability, while across-patient spread is
large. This aligns benchmark evaluation with the cohort-level paradigm of the
clinical literature \cite{TESLA2020,Muller2023}. \textbf{Third}, we quantify the
magnitude gap and synthesize the evidence that the field has systematically set
it aside, providing a reproducible benchmark and protocol that any future
continuous-magnitude predictor can be measured against, while being explicit
about the biological ceiling that bounds what such a predictor can achieve.
